\documentclass{article}

\usepackage{neurips_2025}
\usepackage[utf8]{inputenc}
\usepackage[T1]{fontenc}
\usepackage{hyperref}
\usepackage{url}
\usepackage{booktabs}
\usepackage{amsfonts}
\usepackage{nicefrac}
\usepackage{microtype}
\usepackage{xcolor}
\usepackage{graphicx}
\usepackage{amsmath}

\title{CPU Evaluation of KV Cache Compression Methods on Pythia}

\author{%
  KVPress Reproduction Report\\
  \texttt{EleutherAI/pythia-70m}
}

\begin{document}

\maketitle

\section{Implementation}

The benchmark is implemented in \texttt{pythia\_kvpress\_test.py} and was run as:
\begin{center}
\small
\texttt{python pythia\_kvpress\_test.py ---device cpu}
\end{center}
It evaluates \texttt{EleutherAI/pythia-70m} on CPU in \texttt{float32}. The
dense baseline, \texttt{NoPress}, leaves the KV cache unchanged. The reproduced
methods are \texttt{KnormPress}, which scores tokens from key norms;
\texttt{SnapKVPress}, which scores old tokens by recent-window attention;
\texttt{AdaKVPress}, which applies adaptive head-wise pruning on top of SnapKV
scores; and \texttt{StreamingLLMPress}, which keeps sink tokens and recent
tokens while pruning the middle context.

Perplexity is computed after a 1024-token prefill and on 256 target tokens. For
compressed methods, the prefill runs inside the press hook, so target scoring
uses the compressed cache. Decode throughput is measured only after prefill and
compression; for \texttt{none}, timing starts after prefill only. Therefore the
reported throughput isolates the benefit of a smaller KV cache during decoding,
not the one-time compression cost.

\section{Experimental Setup}

Table~\ref{tab:setup} summarizes the reproduced CPU run. WikiText-2 uses three
samples. PG-19 uses one local sample from \texttt{data/pg19\_sample.txt}, so its
numbers are less stable than the WikiText-2 results.

\begin{table}[!b]
\centering
\caption{Benchmark configuration for \texttt{pythia\_kvpress\_test\_results.csv}.}
\label{tab:setup}
\begin{tabular}{l l}
\toprule
Parameter & Value \\
\midrule
Model & \texttt{EleutherAI/pythia-70m} \\
Device / dtype & CPU / \texttt{float32} \\
Datasets & WikiText-2, PG-19 \\
Context tokens & 1024 \\
Target tokens for PPL & 256 \\
Generated tokens for throughput & 64 \\
Compression ratios & 0.2, 0.4, 0.5, 0.6, 0.8 \\
Warmup windows & 3 \\
Throughput repetitions & 3 \\
SnapKV window size & 64 \\
\bottomrule
\end{tabular}
\end{table}

\section{Results}

\subsection{WikiText-2}

Table~\ref{tab:wikitext} reports perplexity and decode throughput on WikiText-2.
The dense baseline has PPL 32.30 and throughput 65.89 tokens/s. At low
compression ratio 0.2, \texttt{SnapKVPress} and \texttt{StreamingLLMPress}
nearly preserve baseline perplexity while improving throughput by about 40\% and
37\%, respectively. Higher compression usually improves speed but hurts PPL.

\begin{table}[t]
\centering
\caption{WikiText-2 CPU results. PPL is lower better; throughput is higher better.}
\label{tab:wikitext}
\resizebox{\linewidth}{!}{%
\begin{tabular}{lrrrrrrrrrr}
\toprule
Method & \multicolumn{2}{c}{0.2} & \multicolumn{2}{c}{0.4} & \multicolumn{2}{c}{0.5} & \multicolumn{2}{c}{0.6} & \multicolumn{2}{c}{0.8} \\
\cmidrule(lr){2-3}\cmidrule(lr){4-5}\cmidrule(lr){6-7}\cmidrule(lr){8-9}\cmidrule(lr){10-11}
 & PPL & tok/s & PPL & tok/s & PPL & tok/s & PPL & tok/s & PPL & tok/s \\
\midrule
none & 32.30 & 65.89 & 32.30 & 65.89 & 32.30 & 65.89 & 32.30 & 65.89 & 32.30 & 65.89 \\
adakv & 44.55 & 72.27 & 49.80 & 80.31 & 50.95 & 76.64 & 53.17 & 66.54 & 59.90 & 77.75 \\
knorm & 35.43 & 72.18 & 42.94 & 91.44 & 46.69 & 92.96 & 48.06 & 94.05 & 56.43 & 99.22 \\
snapkv & 33.65 & 92.67 & 35.89 & 84.65 & 37.50 & 81.69 & 39.90 & 95.28 & 43.68 & 102.23 \\
streamingllm & 32.71 & 90.56 & 41.40 & 96.81 & 41.86 & 82.15 & 43.10 & 99.29 & 44.77 & 102.97 \\
\bottomrule
\end{tabular}}
\end{table}

\subsection{PG-19}

Table~\ref{tab:pg19} reports the same sweep on PG-19. The dense baseline has
PPL 23.51 and throughput 89.69 tokens/s. \texttt{StreamingLLMPress} is the most
robust method at low compression in this run: at ratio 0.2 it slightly improves
PPL to 23.32 and increases throughput to 94.09 tokens/s. \texttt{SnapKVPress}
also preserves quality well, staying below PPL 25 even at ratio 0.8.

\begin{table}[t]
\centering
\caption{PG-19 CPU results. PPL is lower better; throughput is higher better.}
\label{tab:pg19}
\resizebox{\linewidth}{!}{%
\begin{tabular}{lrrrrrrrrrr}
\toprule
Method & \multicolumn{2}{c}{0.2} & \multicolumn{2}{c}{0.4} & \multicolumn{2}{c}{0.5} & \multicolumn{2}{c}{0.6} & \multicolumn{2}{c}{0.8} \\
\cmidrule(lr){2-3}\cmidrule(lr){4-5}\cmidrule(lr){6-7}\cmidrule(lr){8-9}\cmidrule(lr){10-11}
 & PPL & tok/s & PPL & tok/s & PPL & tok/s & PPL & tok/s & PPL & tok/s \\
\midrule
none & 23.51 & 89.69 & 23.51 & 89.69 & 23.51 & 89.69 & 23.51 & 89.69 & 23.51 & 89.69 \\
adakv & 25.08 & 80.36 & 26.75 & 78.31 & 27.03 & 79.80 & 27.18 & 72.26 & 27.77 & 72.65 \\
knorm & 25.21 & 92.89 & 27.06 & 97.02 & 27.37 & 96.88 & 28.27 & 100.91 & 31.92 & 88.31 \\
snapkv & 23.52 & 91.41 & 23.87 & 98.44 & 23.94 & 100.01 & 23.95 & 97.55 & 24.93 & 98.65 \\
streamingllm & 23.32 & 94.09 & 23.57 & 94.96 & 23.99 & 94.22 & 23.91 & 99.72 & 24.34 & 103.22 \\
\bottomrule
\end{tabular}}
\end{table}

\begin{figure}[t]
\centering
\includegraphics[width=0.9\linewidth]{../results/cpu_base_test/figures/ctxt1024_tar_256_try_cpu_wikitext_tradeoff.png}
\par\medskip
\includegraphics[width=0.9\linewidth]{../results/cpu_base_test/figures/ctxt1024_tar_256_try_cpu_pg19_tradeoff.png}
\caption{Perplexity--throughput tradeoff on WikiText-2 and PG-19. Points closer to the lower-right are better.}
\label{fig:tradeoff}
\end{figure}

\begin{figure}[t]
\centering
\includegraphics[width=0.9\linewidth]{../results/cpu_base_test/figures/ctxt1024_tar_256_try_cpu_wikitext_relative_to_none.png}
\par\medskip
\includegraphics[width=0.9\linewidth]{../results/cpu_base_test/figures/ctxt1024_tar_256_try_cpu_pg19_relative_to_none.png}
\caption{Relative PPL and throughput compared with the no-compression baseline at each compression ratio.}
\label{fig:relative}
\end{figure}

\section{Analysis}

Compression improves decode throughput because each generated token attends to a
shorter KV cache. Since timing starts after compression, the measured speedup
mostly comes from reduced attention work and memory traffic. PPL rises when
pruned tokens contain information needed for the target continuation.

\texttt{SnapKVPress} performs well because recent-window attention is a direct
importance signal for the current context. \texttt{StreamingLLMPress} is also
strong because sink tokens and recent tokens are often enough for stable local
continuation. \texttt{KnormPress} is cheaper but less context-aware, so it loses
more quality at high compression. \texttt{AdaKVPress} does not improve over
SnapKV here; its head-wise budget flexibility adds complexity but does not give
a better tradeoff for this small Pythia-70m CPU run.

The best conservative WikiText-2 point is \texttt{SnapKVPress} at ratio 0.2:
PPL rises from 32.30 to only 33.65, while throughput increases from 65.89 to
92.67 tokens/s. On PG-19, \texttt{StreamingLLMPress} at ratio 0.2 and
\texttt{SnapKVPress} at ratios 0.4--0.6 give the cleanest quality--speed balance.

\section{Conclusion}

The CPU results show the expected tradeoff: cache compression improves
post-compression decode throughput, but aggressive pruning can hurt perplexity.
\texttt{SnapKVPress} is the most consistent quality-preserving method, while
\texttt{StreamingLLMPress} is a strong simple speed baseline. Future runs should
use more PG-19 samples and repeat the sweep on CUDA.

\end{document}
